\documentclass[12pt]{amsart}
%   \overfullrule=5pt
\usepackage[active]{srcltx}
\usepackage{calc,amssymb,amsthm,amsmath,amscd, eucal,ulem}
\usepackage{alltt}
%\usepackage[left=1in,top=1in,right=1in,bottom=1in]{geometry}
\synctex=1
%\usepackage{mathtools}
\RequirePackage[dvipsnames,usenames]{color}
\let\mffont=\sf
\normalem
\input{mabliautoref.sty}
\input{xy}
\xyoption{all}
\input{kmacros3.sty}
\usepackage{tikz}
\usepackage{amsfonts, mathrsfs}
\usepackage{calligra}
\usepackage{stmaryrd} %power series bracket
%\usepackage{dcpic, pictexwd}
%\usepackage{pdfsync}
\usepackage[left=1.1in,top=1in,right=1.1in,bottom=1in]{geometry}
\usepackage{enumitem}

\numberwithin{equation}{theorem}

%\renewcommand{\thefootnote}{\fnsymbol{footnote}}
\newcommand{\D}{\displaystyle}
\newcommand{\til}{\widetilde}
\newcommand{\ol}{\overline}
\newcommand{\F}{\mathbb{F}}
\DeclareMathOperator{\E}{E}
\renewcommand{\:}{\colon}
\newcommand{\eg}{{\itshape e.g.} }
\renewcommand{\m}{\mathfrak{m}}
\renewcommand{\n}{\mathfrak{n}}
\newcommand{\Tt}{{\mathfrak{T}}}
\newcommand{\calC}{\mathcal{C}}
\newcommand{\RamiT}{\mathcal{R_{T}}}
\DeclareMathOperator{\edim}{edim}
\DeclareMathOperator{\length}{length}
\DeclareMathOperator{\monomials}{monomials}
\DeclareMathOperator{\Fitt}{Fitt}
\DeclareMathOperator{\depth}{depth}
\newcommand{\Frob}[2]{{#1}^{1/p^{#2}}}
\newcommand{\Frobp}[2]{{(#1)}^{1/p^{#2}}}
\newcommand{\FrobP}[2]{{\left(#1\right)}^{1/p^{#2}}}
\newcommand{\extends}{extends over $\Tt${}}
\newcommand{\extension}{extension over $\Tt${}}
\newcommand{\fg}{\pi_1^{\textnormal{\'{e}t}}} %Fundamental group
\DeclareMathOperator{\ns}{ns}
%\newcommand{\Spec}{\text{Spec}} %Spectrum
\newcommand{\pSpec}{\text{Spec}^{\circ}} %Puntured Spectrum
\newcommand{\etInCdOne}{\etale{} in codimension $1$}
\DeclareMathOperator{\DIV}{Div}
%\renewcommand{\char}{\charact}


\usepackage{setspace}
\usepackage{hyperref}
%\usepackage{enumerate}
\usepackage{graphicx}
\usepackage[all,cmtip]{xy}

\usepackage{verbatim}

\theoremstyle{theorem}
\newtheorem*{mainthm}{Main Theorem}
\newtheorem*{mainthma}{Main Theorem A}
\newtheorem*{mainthmb}{Main Theorem B}
\newtheorem*{mainthmc}{Main Theorem C}
\newtheorem*{atheorem}{Theorem}


%The todo box!
\def\todo#1{\textcolor{Mahogany}%
{\footnotesize\newline{\color{Mahogany}\fbox{\parbox{\textwidth-15pt}{\textbf{todo:
} #1}}}\newline}}

%What is going on?  Why isn't \O a script O?!?
\renewcommand{\O}{\mathscr O}

\begin{document}
\title{Exercises for Characteristic $p$ commutative algebra\\January 18th, 2017}
\author{Due, January 27th, 2017}
%\address{Department of Mathematics\\ University of Utah\\ Salt Lake City\\ UT 84112}
%\email{schwede@math.utah.edu}


%  \subjclass[2010]{14F18, 13A35, 14B05}

\maketitle

%\section{Reduction to characteristic $p > 0$}
%\chapter{Frobenius and Kunz's theorem}
%\bibliographystyle{skalpha}
%\bibliography{MainBib}

\begin{enumerate}[label=(\arabic*)]
\item Write down a basis for $F^e_* \bF_p[x_1, \ldots, x_n]$ over $\bF_p[x_1, \ldots, x_n]$.
\item Write down a basis for $F^e_* \bF_p(t)[x_1, \ldots, x_n]$ over $\bF_p(t)[x_1, \ldots, x_n]$.
\item If $I = \langle f_1, \ldots, f_t \rangle \subseteq R$ is an ideal in a ring of characteristic $p > 0$, we use $I^{[p^e]}$ to denote the ideal $\langle f_1^{p^e}, \ldots, f_t^{p^e} \rangle$.  Show that $I^{[p^e]}$ is independent of the choice of generators $f_i$ and also show that the formulas:
    \[
    I F^e_* R = F^e_* I^{[p^e]} \text{ and } I R^{1/p^e} = (I^{[p^e]})^{1/p^e}
    \]
    hold.
\item If $R = k[x^2, x^3]$, verify that $R^{1/p}$ is not a free $R$-module for any prime $p$.
\item Use the fact that a Noetherian local ring is regular if and only if it has finite global dimension to prove that if $R$ is a regular local ring, then so is $R_Q$ for any $Q \in \Spec R$.
\item Find the singular locus of the following rings $R$, try to simplify this ideal as much as possible (in particular, take the radical of the ideal defining the singular locus if you can).
\vskip 3pt \noindent
(ie, find an ideal which defines the primes $Q$ in $\Spec R$ such that $R_Q$ is not a regular local ring.  Note you'll have to figure out the dimension of each ring.).
\begin{enumerate}
\item $k = \overline{k}, R = k[x,y,z]/\langle y^3- x^2 z \rangle$.
\item $k = \overline{k}, R = k[x,y,z,w]/\langle z^2 - yw, yz - xw, y^2 - xz \rangle$.
\item (you may want the help of a computer) $k = \overline{k},$  $$R = k[x,y,z,w]/\langle yz-xw,z^3-yw^2,xz^2-y^2w,y^3-x^2z \rangle.$$
\end{enumerate}
\item Consider the following singular local rings defined over $k = \overline{k}, \mathrm{char} k \neq 2$.  See if you can identify which ones remain domains after completion at the identified maximal ideal.
\begin{enumerate}
\item $k[x,y]/\langle y^2 - x^3 \rangle$ localized at the origin $\langle x, y \rangle$.
\item $k[x,y]/\langle y^2 - x^2(x-1) \rangle$ localized at the origin $\langle x,y \rangle$.
\item $k[x,y,z]/\langle xy^2 - z^2\rangle$ localized at the origin $\langle x, y, z \rangle$.
\end{enumerate}
\end{enumerate}



\end{document}


